\section{考频分析}

\subsection{近 5 年题型分值统计}

\begin{center}
\begin{tabular}{c|c|c|c|c|c}
\textbf{年份} & \textbf{信息类} & \textbf{文学类} & \textbf{文言文} & \textbf{古诗} & \textbf{默写} \\
 & \textbf{阅读} & \textbf{阅读} & \textbf{阅读} & \textbf{鉴赏} &  \\ \hline
2024 & 19 分 & 16 分 & 20 分 & 9 分 & 6 分 \\
2023 & 17 分 & 18 分 & 20 分 & 9 分 & 6 分 \\
2022 & 17 分 & 18 分 & 20 分 & 9 分 & 6 分 \\
2021 & 19 分 & 16 分 & 20 分 & 9 分 & 6 分 \\
2020 & 19 分 & 16 分 & 20 分 & 9 分 & 6 分
\end{tabular}
\end{center}

\begin{center}
\begin{tabular}{c|c|c|c|c}
\textbf{年份} & \textbf{语用} & \textbf{写作} & \textbf{总分} & \textbf{选择题占比} \\ \hline
2024 & 20 分 & 60 分 & 150 分 & $\approx 35\%$ \\
2023 & 20 分 & 60 分 & 150 分 & $\approx 33\%$ \\
2022 & 20 分 & 60 分 & 150 分 & $\approx 33\%$ \\
2021 & 20 分 & 60 分 & 150 分 & $\approx 35\%$ \\
2020 & 20 分 & 60 分 & 150 分 & $\approx 35\%$
\end{tabular}
\end{center}

\subsection{各板块命题趋势}

\subsubsection{信息类文本阅读}

\begin{itemize}
  \item \textbf{选材}：科技前沿（人工智能、航天、生命科学）与社科论述（文化传承、美学、社会治理）交替出现
  \item \textbf{新趋势}：2024 年起增加第 4 题多选题、第 5 题开放性简答
  \item \textbf{能力要求}：筛选整合 $\to$ 推理判断 $\to$ 观点评价，层级递进
\end{itemize}

\subsubsection{文学类文本阅读}

\begin{itemize}
  \item \textbf{文体}：散文与小说并重，近年散文考查增多；2024 年出现报告文学
  \item \textbf{主题}：乡土情怀、革命文化、时代精神是三大高频主题
  \item \textbf{高频考点}：句段作用（7 次/5 年）、艺术手法（6 次/5 年）、标题含义（4 次/5 年）
\end{itemize}

\subsubsection{文言文阅读}

\begin{itemize}
  \item \textbf{选材}：《通鉴纪事本末》等纪事本末体持续走热，兼有《说苑》等杂史
  \item \textbf{断句}：2023 年起由客观选择改为自主断句
  \item \textbf{文化常识}：结合语境辨析，不再单独考查死记硬背
\end{itemize}

\subsubsection{古代诗歌鉴赏}

\begin{itemize}
  \item \textbf{朝代}：唐诗（尤其晚唐）\textgreater 宋词 \textgreater 明诗（新增方向）
  \item \textbf{题型}：比较阅读出现频率增加（2 次/5 年）
  \item \textbf{情感类型}：怀古伤今（35\%）、羁旅思乡（25\%）、山水田园（20\%）、酬赠送别（20\%）
\end{itemize}

\subsubsection{语言文字运用}

\begin{itemize}
  \item \textbf{新题型}：成语填空（近义辨析）、语病修改（找错误项）、语句复位（衔接连贯）、修辞赏析
  \item \textbf{命题特点}：语境化、综合化，一材多题趋势明显
\end{itemize}

\subsubsection{写作}

\begin{itemize}
  \item \textbf{命题形式}：材料作文一统天下，近年偏重思辨关系型（对立统一）
  \item \textbf{高频主题}：时代与青年、科技与人文、传承与创新、个人与集体、奋斗与担当
  \item \textbf{评分趋势}：看重思辨深度 \textgreater 文采辞藻；强调结合现实，反对空泛议论
\end{itemize}

\subsection{复习优先级}

\begin{center}
\begin{tabular}{c|c|c}
\textbf{优先级} & \textbf{模块} & \textbf{策略} \\ \hline
\ding{182} 最高 & 写作 & 每周 1 篇限时训练 + 素材积累 \\
\ding{183} & 文言文 & 每日 1 篇精读 + 实词虚词积累 \\
\ding{184} & 信息类阅读 & 掌握定位技巧，限时训练选择题 \\
\ding{185} & 文学类阅读 & 模板化答题 + 典型题型训练 \\
\ding{186} & 诗歌鉴赏 & 分类归纳手法与情感 \\
\ding{187} & 语用 & 新题型专项突破 \\
\ding{188} & 默写 & 考前反复默写易错字 \\
\end{tabular}
\end{center}
